\section{Coefficients and Bott multiplication}
\label{sec:coefficients}

Our first goal is to control multiplication by a chosen coefficient
class, including its interaction with localization. These operations
have two uses. In the arithmetic construction they carry surjectivity
of a valuation map to the degree needed to remove the integral lifting
obstruction. In the geometric argument they translate a relation among
supported coefficient classes into ordinary tame symbols.

We write $G(X)$ for the $K$-theory spectrum of coherent modules on a
Noetherian scheme $X$.  For the regular schemes used below, the canonical
resolution equivalence identifies $K(X)$ with $G(X)$: these schemes have
finite dimension and are affine or quasiprojective over an affine base.
We use Quillen
localization and d\'evissage, including localization by a multiplicative
set and the resulting residue maps
\cite[\S4, Theorem~3 and Corollary~2; \S5, Theorems~4--5;
\S7, Proposition~3.2 and paragraph~3.4]{Quillen1973Higher}.

Let $M=\mathbb S/5$ be the cofiber of multiplication by $5$ on the sphere
spectrum.  Our coefficient convention is
\[
 K_i(X;5)=\pi_i(K(X)\wedge M),\qquad
 G_i(X;5)=\pi_i(G(X)\wedge M).
\]
For $i\geq1$ there is a natural exact sequence
\begin{equation}
\label{eq:coefficient-sequence}
 0\longrightarrow K_i(X)/5\longrightarrow K_i(X;5)
 \xrightarrow{\partial_{\mathrm{coef}}}K_{i-1}(X)[5]
 \longrightarrow0,
\end{equation}
and likewise for $G$. Here $A/5=A/5A$ and $A[5]=\ker(5:A\to A)$ for any abelian group $A$. If $H$ is a field, the identities $K_1(H)=H^\times$ and $K_0(H)=\Z$ give
\begin{equation}
\label{eq:low-coefficients}
 K_1(H;5)=H^\times/H^{\times5},\qquad K_0(H;5)=\mathbb Z/5.
\end{equation}
For a unit $v$, the notation $[v]_5$ denotes its image in $K_1(-;5)$.
We use additive notation in $K$-groups, including groups of unit classes.
For a scheme $X$, the notation $X^{(j)}$ denotes its points of codimension $j$.

We make the coefficient product explicit so that localization acts on
classes with nonzero coefficient boundary as well as on reductions of
integral classes. At the prime five this agrees with the product in
\cite[Proposition~16.1 and Theorem~16.2]{FriedlanderSuslin2002}.

\begin{lemma}[Coefficient products]
\label{lem:moore-product}
The spectrum $M$ has a unique unital pairing
$\mu:M\wedge M\to M$ in the homotopy category.  Multiplication by $5$
on $M$ is nullhomotopic, so all the coefficient groups above are
$\mathbb F_5$-vector spaces.  Tensor product, followed by $\mu$, defines
natural bilinear products in coefficient $K$-theory and the corresponding
module actions on coefficient $G$-theory.
\end{lemma}

\begin{proof}
The cofiber sequence for $M$, together with
$\pi_1\mathbb S=\pi_2\mathbb S=\mathbb Z/2$, gives
$\pi_1M=\pi_2M=0$.  Applying maps into $M$ to this cofiber sequence shows
that restriction along the unit $\eta:\mathbb S\to M$ is injective on
\[
 [M,M]\longrightarrow[\mathbb S,M]=\mathbb Z/5.
\]
Thus $5\operatorname{id}_M=0$.  The cofiber sequence
\[
 M\xrightarrow{5}M\xrightarrow{\eta\wedge1}M\wedge M
 \longrightarrow\Sigma M
\]
therefore permits an extension of $\operatorname{id}_M$ to a pairing
$\mu$.  Its restriction along $1\wedge\eta$ is also the identity:
these endomorphisms of $M$ agree after restriction along $\eta$.
Moreover, $[\Sigma M,M]=0$, by the same cofiber sequence and the two
vanishing homotopy groups.  The extension is consequently unique.
Smashing $5\operatorname{id}_M=0$ with any spectrum proves the assertion
about coefficients.  The products are induced by the usual tensor
pairings of $K$-theory spectra and this fixed pairing on $M$.
\end{proof}

For a base ring $R$ and a class $\beta\in K_2(R;5)$, let
\[
 b: G_i(X;5)\longrightarrow G_{i+2}(X;5)
\]
be left multiplication by the restriction of $\beta$, whenever $X$ is
an $R$-scheme.  We use the same notation for $K$-groups, and write
$b^2=b\circ b$.  In particular $b^2(x)=\beta\cdot(\beta\cdot x)$, with these parentheses. We use this fixed two-factor pairing without assuming an associative coefficient ring structure.

\begin{lemma}[Compatibility with localization]
\label{lem:bott-localization}
The operation $b$ commutes with the maps and boundary homomorphisms in
the $G$-theory localization sequences over $R$, including localization
at a generic point and arbitrary multiplicative localization in the Noetherian affine situations below. Under d\'evissage its action on a residue-field
term is multiplication by the restricted class $\beta$.  It is also
compatible with finite coherent pushforward and flat pullback over $R$.
\end{lemma}

\begin{proof}
Tensoring with a finitely generated projective $R$-module is exact on
coherent modules and preserves every specified support.  Tensor product
therefore gives compatible pairings by $K(R)$ on the supported, whole,
and localized categories.  These pairings respect the null composite
in the localization diagram and hence its cofiber sequence.  On a
d\'evissage category, tensoring is tensoring with the restriction to the
residue field.  The projection formula for a finite map and the usual
flat base-change identity are identities of these tensor functors.
For each term $B$ of this localization triangle, use the pairing
\[
 (K(R)\wedge M)\wedge(B\wedge M)
 \longrightarrow K(R)\wedge B\wedge M\wedge M
 \longrightarrow B\wedge M,
\]
where the last arrow uses the integral tensor action and the fixed
pairing on $M$. These pairings commute with all arrows of the triangle.
Precomposing with the representative
$\beta:\mathbb S^2\to K(R)\wedge M$ gives the asserted diagrams for
$b$; no integral lift of $\beta$ is required.  Its degree is even, so moving it across a boundary contributes
no sign. Passage to a generic point follows by filtered localization. These assertions concern localization boundaries; the coefficient connecting map $\partial_{\mathrm{coef}}$ is a different map.
\end{proof}

We will use the following precise instances of field Bott periodicity.
They concern multiplication by the specified coefficient class, rather
than an unspecified isomorphism between the groups.

\begin{lemma}[The field Bott maps]
\label{lem:bott}
Let $H$ be a perfect field of characteristic different from $5$
containing a primitive fifth root $\zeta_H$.  Suppose that
$\beta_H\in K_2(H;5)$ satisfies
$\partial_{\mathrm{coef}}\beta_H=[\zeta_H]$, and let $b$ denote its
multiplication operation.
\begin{enumerate}
 \item If $\operatorname{cd}_5H\leq3$, then
 \[
 b:K_{2r}(H;5)\xrightarrow{\ \sim\ }K_{2r+2}(H;5)
 \qquad(r\geq1).
 \]
 \item If $\operatorname{cd}_5H\leq2$, then
 \[
 b:K_{2r+1}(H;5)\xrightarrow{\ \sim\ }K_{2r+3}(H;5)
 \qquad(r\geq0).
 \]
 \item If $H$ is finite, then $b:K_i(H;5)\to K_{i+2}(H;5)$
 is an isomorphism for every $i\geq0$.
\end{enumerate}
Over an algebraically closed field of characteristic different from
$5$, the group in degree $2j\geq0$ is generated by $b^j(1)$, and the
odd-degree groups vanish. Here $1\in K_0(H;5)$ is the unit class,
$b^0(1)=1$, and $b^j$ denotes iterated left multiplication.
\end{lemma}

The arithmetic and geometric arguments use the following instances of
the lemma. For number fields and finite extensions of $\mathbb Q_5$,
the $5$-cohomological dimension is two; see
\cite[VI.\S7 and VI.8.2]{Weibel2013}. For a one-variable function field
over a finite extension of $\mathbb Q_5$, the cohomological-dimension
bound \cite[Lemma~59.95.4, Tag~\texttt{0F0T}]{StacksProject} gives
$5$-cohomological dimension at most $2+1=3$. In particular, the actual
multiplication maps are
\begin{equation}
\label{eq:bott-maps-used}
 \begin{aligned}
  b^2 &:K_1(H;5)\xrightarrow{\ \sim\ }K_5(H;5)
    &&\text{for the number and local fields used below},\\
  b^2 &:K_2(J;5)\xrightarrow{\ \sim\ }K_6(J;5)
    &&\text{for the function field of the generic curve}.
 \end{aligned}
\end{equation}
In each case the primitive root and the Bott class are obtained by
restriction from the indicated base. For a finite field $H$ satisfying
the lemma's hypotheses, its finite-field clause also gives
$b^2:K_0(H;5)\xrightarrow{\sim}K_4(H;5)$. In the arithmetic construction
this will raise the target of the valuation homomorphism from degree
zero to degree four.

\begin{proof}[Proof of Lemma~\ref{lem:bott}]
We use the corrected, strongly convergent motivic spectral sequence
with its finite-coefficient multiplicative structure:
\begin{equation}
\label{eq:motivic-ss}
 E_2^{p,q}=H^{p-q}(H,\mathbb Z/5(-q))
 \Longrightarrow K_{-p-q}(H;5).
\end{equation}
This form of the spectral sequence is recorded in
\cite[VI.4.2, Addendum~VI.4.2.1, and VI.4.4]{Weibel2013}.
For the corrected construction, see
\cite[Introduction and Theorem~6.1]{Suslin2003}
(the locators use the \emph{Trudy} pagination).
The product is the usual $K$-theory product on the abutment and the
motivic cup product on $E_2$
\cite[Proposition~16.1 and Theorem~16.2]{FriedlanderSuslin2002}.
The Moore-space comultiplication used there corresponds by stable
duality to a unital pairing on $M$.  Its coefficient product agrees with
ours by the uniqueness in Lemma~\ref{lem:moore-product}.

Write an $E_2$-term as $H^s(H,\mathbb Z/5(j))$, of total $K$-degree
$2j-s$.  Terms of negative weight vanish. For $j\geq0$, the defining
motivic complex
$C_*\mathbb Z_{\mathrm{tr}}(\mathbb G_m^{\wedge j})[-j]$
is concentrated in cohomological degrees at most $j$. Global sections
on the Nisnevich site of a field are exact, since every covering of
the field point has a splitting, and derived reduction modulo five does not
increase this upper degree. Thus the field motivic group vanishes
when $s>j$. For $s\leq j$, the norm-residue and
Beilinson--Lichtenbaum comparison theorems identify it with
\begin{equation}
\label{eq:motivic-etale}
 H^s_{\mathrm{et}}(H,\mu_5^{\otimes j})
\end{equation}
\cite[Theorems~6.16--6.17]{Voevodsky2011}.
These are the Rost--Voevodsky norm-residue theorem and its motivic
comparison consequence, with the completion of the norm-residue
argument in \cite{Weibel2009}; the earlier motivic formulation is
\cite{SuslinVoevodsky2000}.
In particular the terms of negative cohomological degree vanish.
Under the dimension bounds below, only finitely many terms contribute
to each total degree, so the resulting filtrations are finite.
In these coordinates a differential $d_\rho$, for $\rho\geq2$, goes
from $(s,j)$ to $(s+2\rho-1,j+\rho-1)$.

We first derive the algebraically closed case from this same spectral
sequence. If $\overline H$ is an algebraic closure of $H$, only the
terms with $s=0$ survive. In each degree $2j\geq0$ there is exactly one
term, $H^0(\overline H,\mu_5^{\otimes j})\cong\F_5$, and in odd
degree there is none. There are no possible differentials. Thus the
edge map is an isomorphism in every even degree, and
\[
 K_{2j}(\overline H;5)\cong\F_5,\qquad
 K_{2j+1}(\overline H;5)=0\qquad(j\geq0).
\]
The coefficient sequence surjects
$K_2(\overline H;5)$ onto $\overline H^\times[5]=\mu_5(\overline H)$.
Both groups have dimension one, so this boundary is an isomorphism.
Consequently the restriction of $\beta_H$ is nonzero, and so is its
edge image. Multiplicativity now shows that each $b^j(1)$ over
$\overline H$, computed by iterated left multiplication, has nonzero
edge image and generates the corresponding even group. This proves
the last assertion of the lemma without a
separate rigidity input.

By naturality, the edge image
$\alpha\in H^0(H,\mu_5)$ of $\beta_H$ is nonzero: restriction on this
$H^0$ group is injective, and its geometric image is the nonzero class
just computed. The edge image of $b^r(1)$ is
$\alpha^{\otimes r}$. Since $H$ contains $\mu_5$, this generates the
one-dimensional group $H^0(H,\mu_5^{\otimes r})$. Hence this whole
line consists of permanent cycles for every $r\geq1$. Only
$\alpha\ne0$ is needed; no choice of an identification of $\alpha$
with the specified primitive root is required.

Suppose first that $\operatorname{cd}_5H\leq3$.  In total degree
$2r$, the only possible terms are
\[
 H^0(H,\mu_5^{\otimes r}),\qquad
 H^2(H,\mu_5^{\otimes(r+1)}).
\]
There is no incoming differential to either term.  The only possible
outgoing differential from the first is $d_2$ to an $H^3$-term, and
it vanishes by permanence of the Bott power.  Every outgoing
differential from the second has cohomological degree at least $5$
and is zero.  We obtain a two-step filtration, expressed as
\begin{equation}
\label{eq:even-field-filtration}
 0\longrightarrow H^2(H,\mu_5^{\otimes(r+1)})
 \longrightarrow K_{2r}(H;5)
 \longrightarrow H^0(H,\mu_5^{\otimes r})
 \longrightarrow0.
\end{equation}
Multiplication by $\beta_H$ preserves these filtrations and acts on
both graded pieces by tensoring with $\alpha$.  Each such map is an
isomorphism, proving (1).

If $\operatorname{cd}_5H\leq2$, the sole term in total degree
$2r+1$ is $H^1(H,\mu_5^{\otimes(r+1)})$, with no possible nonzero
differential.  Multiplication by $\beta_H$ again acts by tensoring
with $\alpha$.  This proves (2).  For a finite field, the same
calculation applies in every positive degree since its
$5$-cohomological dimension is one.  In degree zero, the class $1$
generates $K_0(H;5)$ and $b(1)=\beta_H$ generates
$K_2(H;5)=H^0(H,\mu_5)$.  This proves (3).
\end{proof}
